\subsection{Описание функций}

\subsection*{Функция для создания графа}

\begin{verbatim}
    vector<vector<pair<int, int>>> generateGraph(int n, int m, int q, int r) {
        vector<vector<pair<int, int>>> adj_list(n);
        std::random_device rd;
        std::mt19937 gen(rd());
        std::uniform_int_distribution<> dist(q, r);
        vector<vector<bool>> was(n, vector<bool>(n, false));
    
        for (int edge_count = 0; edge_count < min(n * (n - 1) / 2, m);) {
            int i = rand() % n, j = rand() % n;
            if (i != j && !was[i][j]) {
                int weight = dist(gen);
                adj_list[i].push_back({ j, weight });
                adj_list[j].push_back({ i, weight });
                was[i][j] = was[j][i] = true;
                ++edge_count;
            }
        }
        return adj_list;
    }
    \end{verbatim}

Эта функция \texttt{generateGraph} предназначена для создания случайного взвешенного графа. Она создает список смежности для графа с заданным количеством узлов и случайным набором рёбер, каждое из которых имеет случайный вес в заданном диапазоне.

\subsubsection*{Назначение:}
Генерация случайного неориентированного графа с определённым числом узлов и рёбер, где веса рёбер выбираются случайным образом в указанном диапазоне. Такой граф может использоваться для тестирования алгоритмов на больших случайных данных.

\subsubsection*{Параметры:}
\begin{itemize}
    \item \texttt{int n} — количество узлов (вершин) графа.
    \item \texttt{int m} — максимальное количество рёбер.
    \item \texttt{int q} — минимально возможный вес рёбер.
    \item \texttt{int r} — максимально возможный вес рёбер.
\end{itemize}

\subsubsection*{Возвращаемое значение:}
Возвращает граф в виде списка смежности, представленного вектором векторов пар. Каждая пара в списке смежности содержит узел, с которым соединена вершина, и вес рёбра.



\subsection*{Алгоритм Дейкстры}

\begin{verbatim}
    double dijkstra(vector<int>& dist, vector<int>& up, 
                    vector<vector<pair<int, int>>> graph, int n, int d, int start) {
        vector<bool> processed(n, false);
        fill(dist.begin(), dist.end(), INT_MAX);
        fill(up.begin(), up.end(), -1);
        dist[start] = 0;
        DHeap dHeap(d);
        dHeap.Insert(start, 0);
    
        auto begin = chrono::high_resolution_clock::now();
    
        while (!dHeap.IsEmpty()) {
            auto [i, currentDist] = dHeap.ExtractMin();
            if (processed[i]) continue;
            processed[i] = true;
            
            for (auto& [j, weight] : graph[i]) {
                int newDist = currentDist + weight;
                if (newDist < dist[j]) {
                    dist[j] = newDist;
                    up[j] = i;
                    dHeap.Insert(j, newDist);
                }
            }
        }
    
        auto end = chrono::high_resolution_clock::now();
        return chrono::duration<double>(end - begin).count();
    }
    \end{verbatim}

Функция \texttt{dijkstra} реализует алгоритм Дейкстры для нахождения кратчайших путей от заданной вершины до всех остальных. Алгоритм использует специальную структуру данных \texttt{DHeap}, которая позволяет ускорить поиск на графах.

\subsubsection*{Назначение:}
Вычисление кратчайших расстояний от начального узла до всех других узлов графа с неотрицательными весами рёбер.

\subsubsection*{Параметры:}
\begin{itemize}
    \item \texttt{vector<int>\& dist} — вектор, в который записываются минимальные расстояния от начального узла до всех других узлов.
    \item \texttt{vector<int>\& up} — вектор предшественников, где для каждого узла сохраняется его предыдущий узел в кратчайшем пути.
    \item \texttt{vector<vector<pair<int, int>>> graph} — граф, представленный в виде списка смежности.
    \item \texttt{int n} — количество узлов в графе.
    \item \texttt{int d} — степень структуры данных \texttt{DHeap}, которая определяет количество детей в каждом узле кучи.
    \item \texttt{int start} — начальный узел, с которого начинается поиск кратчайших путей.
\end{itemize}

\subsubsection*{Возвращаемое значение:}
Возвращает время выполнения алгоритма в секундах, что позволяет оценить его эффективность на данном графе.
